\documentclass{article}
\usepackage{iclr2027_conference,times}
\input{math_commands.tex}
\usepackage{amsmath,amssymb,booktabs,microtype,array}
\usepackage{fontspec}
\usepackage{hyperref}
\usepackage{url}

\title{Đề xuất nghiên cứu: Characteristic Spectral Transform\\
cho chưng cất tri thức mô hình ngôn ngữ}

\author{Bản đề xuất kỹ thuật bằng tiếng Việt}

% Bật dòng dưới nếu cần bản không ẩn tác giả.
% \iclrfinalcopy

\newcommand{\cst}{\textsc{CST}}
\newcommand{\sg}{\operatorname{stopgrad}}

\begin{document}
\maketitle

\begin{abstract}
Đề xuất này trình bày Characteristic Spectral Transform (\cst), một hàm mất
mát phụ trợ dùng để truyền hình học biểu diễn từ mô hình giáo viên sang mô hình
học sinh.  Thay vì ép hai hidden state bằng nhau theo từng tọa độ, học một
projector đổi chiều, hoặc so trực tiếp các trị riêng, \cst{} mô tả hình học của
hidden state bằng đường cong
$\Phi_H(\gamma)=\log\det(I+\gamma G_H)$ tại nhiều giá trị $\gamma>0$.
$G_H$ là ma trận Gram đã trừ trung bình và chuẩn hóa năng lượng.  Vì vậy, hai
mô hình có hidden width khác nhau vẫn so sánh được bằng các số vô hướng.
Trong đường train, log-determinant được tính bằng Cholesky theo batch, không sử
dụng SVD hay eigendecomposition.  Tài liệu giải thích từng công thức, thuật toán,
độ phức tạp, cấu hình thực nghiệm và cách diễn giải kết quả ban đầu.
\end{abstract}

\section{Mục tiêu và câu hỏi nghiên cứu}

Chưng cất tri thức (knowledge distillation, KD) thường yêu cầu mô hình học sinh
khớp phân phối đầu ra của mô hình giáo viên.  Với một token, nếu $z_s,z_t$ là
logit của student và teacher, một dạng loss đầu ra là
\begin{equation}
 \mathcal L_{\mathrm{KD}}
 =D\!\left(
   \operatorname{softmax}(z_s/T),
   \operatorname{softmax}(z_t/T)
 \right),
 \label{eq:output-kd}
\end{equation}
trong đó $D$ có thể là KL, skew-KL hoặc một divergence khác và $T$ là nhiệt độ.
Loss này trả lời câu hỏi ``student có dự đoán giống teacher không?'', nhưng
không trực tiếp trả lời câu hỏi ``student có tổ chức biểu diễn trung gian giống
teacher không?''.

Đề xuất đặt ra ba câu hỏi:
\begin{enumerate}
 \item Có thể truyền cấu trúc phổ của hidden state mà không cần hai hidden width
 bằng nhau hay không?
 \item Có thể tránh projector học được, teacher centroids, SVD và
 eigendecomposition trong đường train hay không?
 \item Có thể làm điều đó với overhead đủ nhỏ cho on-policy KD hay không?
\end{enumerate}

Mục tiêu cuối cùng là
\begin{equation}
 \boxed{
 \mathcal L_{\mathrm{total}}
 =\mathcal L_{\mathrm{KD}}
 +\lambda_{\cst}\mathcal L_{\cst}}
 \label{eq:total}
\end{equation}
và giữ nguyên toàn bộ loss KD đầu ra hiện có.

\section{Ký hiệu và dữ liệu đầu vào}

Xét một minibatch sau khi student đã sinh response.  Tại một layer, gom hidden
state của các response token hợp lệ thành
\begin{equation}
 H\in\R^{m\times d}.
\end{equation}
Ở đây:
\begin{itemize}
 \item $m$ là số token được chọn; cấu hình mặc định là $m\leq64$;
 \item $d$ là hidden width, ví dụ student và teacher có thể lần lượt là
 $d_s=1536$ và $d_t=5120$;
 \item mỗi hàng $H_{i,:}$ là biểu diễn của một token;
 \item teacher và student dùng đúng cùng vị trí token, nhưng không cần cùng
 số chiều đặc trưng.
\end{itemize}

Mask response lấy trực tiếp từ label:
\begin{equation}
 \mathcal I=\{i:\texttt{label}_i\neq-100\}.
\end{equation}
Nếu $|\mathcal I|>m_{\max}$, ta lấy ngẫu nhiên một tập con một lần trong mỗi
training step.  Tập chỉ số đó được dùng chung cho mọi layer, student và teacher.

\section{Từ hidden state tới hình học chuẩn hóa}

\subsection{Bước 1: trừ trung bình theo token}

Vector trung bình của layer là
\begin{equation}
 \bar h=\frac{1}{m}\sum_{i=1}^{m}H_{i,:}\in\R^d.
\end{equation}
Hidden state sau khi center:
\begin{equation}
 X=H-\mathbf 1\bar h^\top\in\R^{m\times d}.
 \label{eq:center}
\end{equation}
Thao tác này loại bỏ phép tịnh tiến chung.  Nếu mọi token cùng được cộng một
vector $b$, tức $H'=H+\mathbf1b^\top$, thì $X'=X$.

\subsection{Bước 2: tính tổng năng lượng}

\begin{equation}
 e=\lVert X\rVert_F^2
 =\sum_{i=1}^{m}\sum_{k=1}^{d}X_{ik}^2.
 \label{eq:energy}
\end{equation}
Nếu $e\leq\epsilon$, layer gần như hằng theo token và không chứa hình học đáng
tin cậy; implementation bỏ qua layer này để tránh chia cho số gần zero.

\subsection{Bước 3: ma trận Gram token}

\begin{equation}
 G_H=\frac{XX^\top}{e+\epsilon}\in\R^{m\times m}.
 \label{eq:gram}
\end{equation}
Phần tử $(G_H)_{ij}$ đo inner product chuẩn hóa giữa token $i$ và token $j$.
Ma trận này là positive semidefinite (PSD), và khi $e\gg\epsilon$:
\begin{equation}
 \operatorname{Tr}(G_H)
 =\frac{\operatorname{Tr}(XX^\top)}{e+\epsilon}
 \approx1.
\end{equation}
Do đó $G_H$ giữ lại \emph{cách năng lượng được phân bố giữa các hướng}, nhưng
loại bỏ độ lớn tổng thể.

Ba bất biến quan trọng là
\begin{align}
 G_{H+\mathbf1b^\top}&=G_H &&\text{(tịnh tiến)},\\
 G_{aH}&=G_H,\quad a\neq0 &&\text{(scale đẳng hướng)},\\
 G_{HQ}&=G_H,\quad QQ^\top=I &&\text{(xoay hệ tọa độ)}.
\end{align}
Đây là lý do \cst{} không cần học ánh xạ tọa độ từ student sang teacher.

\section{Characteristic Spectral Transform}

\subsection{Định nghĩa trung tâm}

Với một scale $\gamma>0$, định nghĩa
\begin{equation}
 \boxed{
 \Phi_H(\gamma)=\log\det(I+\gamma G_H)}.
 \label{eq:cst}
\end{equation}
Vì $G_H\succeq0$, mọi trị riêng của $I+\gamma G_H$ lớn hơn hoặc bằng 1;
log-determinant vì vậy hữu hạn và không âm.

Nếu chỉ dùng để phân tích, gọi $p_1,\ldots,p_r$ là các trị riêng dương của
$G_H$.  Khi đó
\begin{equation}
 \Phi_H(\gamma)=\sum_{i=1}^{r}\log(1+\gamma p_i).
 \label{eq:spectrum-view}
\end{equation}
Công thức~\ref{eq:spectrum-view} giải thích ý nghĩa phổ, nhưng implementation
train \textbf{không tính} $p_i$.

\subsection{Tại sao cần nhiều giá trị $\gamma$?}

Giả sử hai phổ đều có tổng bằng 1:
\begin{align}
 p^{(A)}&=(0.5,0.5),\\
 p^{(B)}&=(0.9,0.1).
\end{align}
Ở scale nhỏ, khai triển Taylor cho ta
\begin{equation}
 \Phi_H(\gamma)
 =\gamma\operatorname{Tr}(G_H)
 -\frac{\gamma^2}{2}\operatorname{Tr}(G_H^2)
 +O(\gamma^3).
 \label{eq:taylor}
\end{equation}
Vì trace gần 1, số hạng bậc nhất của mọi representation gần giống nhau.  Sự
khác biệt bắt đầu rõ ở
\begin{equation}
 \operatorname{Tr}(G_H^2)=\sum_i p_i^2,
\end{equation}
đại lượng đo mức tập trung phổ.  Với ví dụ trên, giá trị này lần lượt là 0.5 và
0.82.  Scale lớn hơn tiếp tục cung cấp thông tin về các moment bậc cao và số
hướng hoạt động.  Vì vậy một $gamma$ duy nhất có thể chỉ nhìn thấy một phần
hình học; nhiều $gamma$ mô tả cả một đường cong đặc trưng.

\subsection{Tính đặc trưng của transform}

Đạo hàm theo scale là
\begin{equation}
 \Phi'_H(\gamma)=\sum_{i=1}^{r}\frac{p_i}{1+\gamma p_i}.
 \label{eq:derivative}
\end{equation}
Vế phải là một hàm hữu tỉ; các pole tại $-1/p_i$ xác định các trị riêng dương.
Do đó, với phổ hữu hạn, nếu hai transform bằng nhau trên một khoảng $gamma$
mở thì hai normalized spectrum bằng nhau, kể cả multiplicity.  Trong train ta
không quan sát cả khoảng liên tục mà dùng Monte Carlo tại các scale mới qua mỗi
step.

\section{Loss \cst{} nhiều scale và nhiều layer}

Lấy log-uniform để mọi bậc độ lớn có xác suất như nhau:
\begin{align}
 u_j&\sim\mathcal U(\log\gamma_{\min},\log\gamma_{\max}),\\
 \gamma_j&=\exp(u_j),\qquad j=1,\ldots,q.
 \label{eq:gamma}
\end{align}
Student và teacher bắt buộc dùng cùng $\gamma_j$.  Loss tại cặp layer $ell$:
\begin{equation}
 \mathcal L_{\cst}^{(\ell)}
 =\frac1q\sum_{j=1}^{q}
 \left(
  \Phi_{H_s^{(\ell_s)}}(\gamma_j)
  -\sg\left[\Phi_{H_t^{(\ell_t)}}(\gamma_j)\right]
 \right)^2.
 \label{eq:layerloss}
\end{equation}
Với $n$ cặp layer:
\begin{equation}
 \boxed{
 \mathcal L_{\cst}=\frac1n\sum_{\ell=1}^{n}
 \mathcal L_{\cst}^{(\ell)}}.
 \label{eq:multilayer}
\end{equation}

Ký hiệu stop-gradient có hai ý nghĩa thực tế: teacher chạy dưới
\texttt{no\_grad}, và gradient của Equation~\ref{eq:total} chỉ cập nhật
student.  CST cũng không sinh thêm module nào ở inference.

\section{Tính log-determinant mà không eig/SVD}

\subsection{Chọn không gian determinant nhỏ hơn}

Thông thường $m=64\ll d$, nên dùng Gram token $XX^\top$ kích thước $64\times64$.
Nếu có trường hợp $d<m$, định lý determinant Sylvester cho phép
\begin{equation}
 \det(I_m+\gamma XX^\top/e)
 =\det(I_d+\gamma X^\top X/e).
 \label{eq:sylvester}
\end{equation}
Vì vậy luôn chọn
\begin{equation}
 r=\min(m,d),\qquad B\in\R^{r\times r}.
\end{equation}

\subsection{Dùng một Gram cho mọi gamma}

Mỗi layer chỉ tính một lần
\begin{equation}
 B=\frac{XX^\top}{e+\epsilon}
 \quad\text{hoặc}\quad
 B=\frac{X^\top X}{e+\epsilon}.
\end{equation}
Sau đó broadcast toàn bộ $q$ scale:
\begin{equation}
 A=I_{[1,r,r]}+gamma_{[q,1,1]}B_{[1,r,r]}
 \in\R^{q\times r\times r}.
 \label{eq:batchA}
\end{equation}
Không được tính lại $X X^\top$ bên trong vòng lặp gamma.

\subsection{Batched Cholesky}

Với từng $A_j\succ0$:
\begin{align}
 A_j&=L_jL_j^\top,\\
 \log\det(A_j)
 &=2\sum_{k=1}^{r}\log(L_j)_{kk}.
 \label{eq:cholesky}
\end{align}
Một lệnh \texttt{torch.linalg.cholesky\_ex(A)} xử lý cả batch $q$.  Mọi phép
tính phổ chạy float32 dù forward chính là fp16/bf16.  Gram được symmetrize:
\begin{equation}
 B\leftarrow\frac12(B+B^\top).
\end{equation}
Nếu Cholesky hiếm khi thất bại do sai số số học, retry một lần với
$A_j+10^{-6}I$.

\subsection{Độ phức tạp}

Chi phí mỗi layer xấp xỉ
\begin{equation}
 O(mdr)+O(qr^3),\qquad r=\min(m,d).
\end{equation}
Với $m=r=64$, $q=2$, $n=4$, phần determinant rất nhỏ so với forward teacher
14B.  Log thực nghiệm đo khoảng 2.2ms cho CST kernel, trong khi một optimizer
step đầy đủ của on-policy pipeline khoảng 22.6 giây.  Con số 2.2ms không bao
gồm toàn bộ chi phí giữ hidden state, nên báo cáo cuối cùng vẫn cần đo end-to-end
so với output-KD baseline.

\section{CST khác gì NuNo và RPT?}

\begin{table}[t]
\caption{Phân biệt ba representation objective.}
\label{tab:compare}
\centering
\footnotesize
\setlength{\tabcolsep}{2pt}
\begin{tabular}{>{\raggedright\arraybackslash}p{2.1cm}
                >{\raggedright\arraybackslash}p{3.4cm}
                >{\raggedright\arraybackslash}p{3.4cm}
                >{\raggedright\arraybackslash}p{3.4cm}}
\toprule
 & NuNo & RPT & CST\\
\midrule
Đối tượng & Nuclear norm có centroid & Profile tích lũy trị riêng &
$\log\det(I+\gamma G)$\\
Projector & Có & Không & Không\\
Centroid & Có & Không & Không\\
Eig/SVD khi train & Tránh bằng Newton--Schulz & Phân rã trị riêng & Không\\
Nhiều scale & Không & Không & Có\\
\bottomrule
\end{tabular}
\end{table}

RPT là baseline phân tích hợp lệ nhưng không phải cách triển khai CST.  Không
được đơn giản hóa Equation~\ref{eq:cst} thành loss khớp trực tiếp eigenvalue.
Tương tự, projector đã tồn tại cho NuNo không được tái sử dụng trong CST path.

\section{Cấu hình thực nghiệm đề xuất}

\begin{table}[t]
\caption{Cấu hình Qwen và CST hiện tại.}
\label{tab:config}
\centering
\small
\begin{tabular}{ll}
\toprule
Teacher $\rightarrow$ student & Qwen2.5-14B $\rightarrow$ 1.5B\\
Global batch / epochs & 64 / 2\\
Learning rate / scheduler & $5\times10^{-6}$ / cosine\\
Warmup / max length & 0.1 / 1024\\
Số layer / depth range & 4 / $[0.20,0.85]$\\
Token tối đa $m$ & 64\\
Số gamma $q$ & 2\\
Khoảng gamma & $[10^{-2},10^2]$, log-uniform\\
Khoảng cách transform & squared $\ell_2$\\
$\lambda_{\cst}$ & 0.003\\
\bottomrule
\end{tabular}
\end{table}

Tại đầu run, raw CST loss thường nằm khoảng 30--65, còn output-KD khoảng
1.2--1.4.  Nếu dùng nhầm weight 0.2 của NuNo, thành phần CST có thể lớn khoảng
12 và áp đảo KD.  Với $\lambda_{\cst}=0.003$, contribution điển hình là
\begin{equation}
 0.003\times60=0.18,
\end{equation}
tức khoảng 10--15\% output loss.  Vì hai loss có scale khác nhau, không nên
đồng nhất $\lambda_{\cst}$ với $\lambda_{\mathrm{NNM}}$.

Grid tune đề xuất:
\begin{align}
 \lambda_{\cst}&\in\{0.001,0.003,0.01,0.03\},\\
 m&\in\{32,64,128\},\\
 q&\in\{1,2,4\},\\
 n&\in\{2,4\}.
\end{align}
Chỉ thay một trục trong mỗi ablation và giữ seed, dữ liệu, output-KD, decoding
và evaluator cố định.

\section{Kết quả ban đầu và cách đọc đúng}

Run hiện có dùng 10K training examples, hai epochs, trong khi bảng tham chiếu
NuNo dùng khoảng 78K examples.  Vì ngân sách chưa bằng nhau, đây là kết quả
pilot chứ chưa phải chứng minh CST hơn phương pháp khác.

\begin{table*}[t]
\caption{Accuracy (\%) trên Qwen2.5-14B $\rightarrow$ 1.5B. Các dòng tham
chiếu lấy từ manuscript NuNo; CST dùng 10K examples. MATH trong average dùng
exact-match để gần bảng cũ nhất có thể.}
\label{tab:results}
\centering
\small
\setlength{\tabcolsep}{3pt}
\begin{tabular}{lrrrrrrrrr}
\toprule
Method & GSM8K & GSM+ & MATH & MBPP & SciQ & STEM & Pro-Math & BBH & Avg.\\
\midrule
Student tham chiếu & 57.1 & 38.8 & 16.9 & 38.4 & 85.5 & 49.5 & 34.4 & 36.5 & 44.6\\
NuNo tham chiếu & 65.2 & 46.8 & 32.2 & 44.4 & 86.2 & 51.7 & 42.0 & 44.9 & 51.7\\
\midrule
CST, 10K & 65.58 & 46.79 & 14.92 & 46.80 & 86.90 & 51.95 & 44.49 & 41.87 & 49.91\\
\bottomrule
\end{tabular}
\end{table*}

CST 10K tăng 5.3 điểm so với base student tham chiếu và gần bằng NuNo trên
GSM8K, GSM-Plus, SciQ và STEM.  Nó cao hơn ở MBPP và Pro-Math nhưng thấp hơn ở
BBH.  Cột MATH cần giải thích riêng.

Harness hiện tại trả về hai metric trên cùng 5,000 generation:
\begin{align}
 \texttt{exact\_match}&=14.92,\\
 \texttt{math\_verify}&=45.06.
\end{align}
Trong 4,454 sample file đã kiểm tra, confusion giữa hai metric là:
\begin{center}
\begin{tabular}{ccr}
\toprule
Exact & Verify & Số mẫu\\
\midrule
0 & 0 & 2,372\\
0 & 1 & 1,409\\
1 & 1 & 669\\
1 & 0 & 4\\
\bottomrule
\end{tabular}
\end{center}
Nhiều câu trong nhóm $(0,1)$ có đáp án toán đúng và được đóng khung
\texttt{\textbackslash boxed\{...\}}, nhưng thiếu đúng chuỗi ``Final Answer''
mà exact parser mong đợi.  Vì vậy 14.92 bị ảnh hưởng mạnh bởi format; ngược lại,
không được so thẳng 45.06 với MATH 32.2 của một evaluator version khác.  Cách
so sánh hợp lệ là chạy CST, NuNo, RPT và output-KD bằng cùng container, rồi báo
cả exact và symbolic verification.

\section{Kế hoạch hoàn thiện paper}

\paragraph{Thí nghiệm bắt buộc.}
\begin{enumerate}
 \item Chạy đủ khoảng 78K mẫu cho Qwen, ít nhất ba seed cho cấu hình chính.
 \item Train lại output-KD, NuNo, RPT và CST với cùng data order và budget.
 \item Đóng băng phiên bản Transformers, lm-eval, vLLM, task YAML và prompt.
 \item Chạy cặp Gemma2-9B $\rightarrow$ 2B để kiểm tra khả năng tổng quát.
 \item Báo wall-clock, peak memory và throughput end-to-end, không chỉ CST
 kernel time.
\end{enumerate}

\paragraph{Ablation quan trọng.}
Ngoài $m,q,n,\lambda$, cần so sánh center và không center; log-uniform với fixed
grid; squared $\ell_2$ với smooth-$\ell_1$; shared gamma với gamma khác nhau
(negative control); và shared token indices với sampling độc lập (negative
control).  Hai negative control dự kiến kém hơn vì phá tính so sánh trực tiếp.

\paragraph{Phân tích representation.}
Eigendecomposition chỉ được dùng offline để vẽ effective rank, entropy và rank
profile theo layer.  Nó là công cụ kiểm tra CST có thực sự thay đổi geometry
hay không, không được đưa trở lại training objective.

\section{Rủi ro và giới hạn}

\begin{itemize}
 \item Hai gamma mỗi step chỉ là estimator thưa của một đường cong liên tục.
 \item Global token sampling có thể trộn các sequence dài/ngắn không đều.
 \item Loss scale phụ thuộc khoảng gamma và cần tune lại khi đổi cấu hình.
 \item Kết quả hiện tại mới có một seed, một model pair và 10K examples.
 \item Benchmark accuracy không đo safety; distillation có thể sao chép bias và
 lỗi của teacher.
\end{itemize}

\section{Kết luận đề xuất}

\cst{} biến bài toán khớp geometry thành bài toán khớp một hàm vô hướng nhiều
scale.  Trọng tâm phương pháp không phải eigenvalue matching, mà là objective
transform-level
\begin{equation*}
 \Phi_H(\gamma)=\log\det(I+\gamma G_H).
\end{equation*}
Thiết kế này giải quyết tự nhiên hidden-width mismatch, loại bỏ projector và
teacher centroid, đồng thời có đường tính toán Cholesky nhỏ và ổn định.  Kết quả
10K ban đầu đủ để biện minh cho full-budget study, nhưng mọi tuyên bố vượt NuNo
phải chờ controlled reproduction với evaluator cố định.

\appendix
\section{Pseudocode gần với implementation}
\begin{verbatim}
idx = response_mask.nonzero()
idx = random_subset(idx, max_tokens=64)  # một lần mỗi step
gamma = exp(uniform(log(1e-2), log(1e2), size=[2]))

losses = []
for s_layer, t_layer in layer_pairs:
    Xs = center(student_hidden[s_layer][idx]).float()
    Xt = center(teacher_hidden[t_layer][idx]).float().detach()
    Bs = smaller_gram(Xs) / (frobenius_sq(Xs) + eps)
    Bt = smaller_gram(Xt) / (frobenius_sq(Xt) + eps)
    As = I[None] + gamma[:, None, None] * Bs[None]
    At = I[None] + gamma[:, None, None] * Bt[None]
    phi_s = 2 * log(diag(cholesky(As))).sum(-1)
    phi_t = 2 * log(diag(cholesky(At))).sum(-1)
    losses.append(mean((phi_s - stopgrad(phi_t)) ** 2))

L_cst = mean(losses)
L_total = L_output_kd + lambda_cst * L_cst
\end{verbatim}

\section{Checklist kiểm tra implementation}
\begin{itemize}
 \item Không có projector trong CST path.
 \item Không có centroid, K-means hay random projection.
 \item Không gọi SVD/eigvalsh trong main loss.
 \item Teacher và student dùng cùng token indices, layer mapping và gamma.
 \item Teacher transform stop-gradient.
 \item Gram chỉ tính một lần mỗi layer rồi dùng lại cho mọi gamma.
 \item Gamma được batch trong một Cholesky call.
 \item Dùng determinant space nhỏ hơn và float32.
 \item Degenerate layer được skip an toàn.
 \item Unit tests kiểm tra invariance, Sylvester, slogdet reference, gradient,
 teacher stop-gradient và projector-free dispatch.
\end{itemize}

\end{document}
